\documentclass[aspectratio=169,11pt]{beamer}
\usepackage{../../../shared/theme/esmad_beamer_theme}
\usepackage{amsmath,amssymb}
\usepackage{booktabs}

\title[Transportability]{Transportability and External Validity}
\subtitle{Moving Causal Evidence Across Populations, Sites, and Time}
\date{\today}

\newcommand{\E}{\mathbb{E}}

\begin{document}

\begin{frame}
\titlepage
\end{frame}

\begin{frame}{Learning Goals}
\begin{itemize}
\item distinguish internal validity from external validity;
\item define study populations, target populations, and target estimands;
\item explain effect modification and why it drives transportability;
\item use standardization and weighting to generalize effects;
\item diagnose support failures and covariate shift;
\item reason about site heterogeneity, temporal drift, and deployment differences.
\end{itemize}
\end{frame}

\begin{frame}{Outline}
\tableofcontents
\end{frame}

\section{Internal and External Validity}

\begin{frame}{Internal Validity}
Internal validity asks whether the effect is identified in the study population.

A study can be internally valid when:
\begin{itemize}
\item treatment assignment is well understood;
\item confounding is controlled or randomized;
\item measurement is reliable;
\item interference is handled or absent.
\end{itemize}
\end{frame}

\begin{frame}{External Validity}
External validity asks whether the study effect informs another population, site, or time period.

\begin{alertblock}{Key point}
A precisely estimated internally valid effect can still be irrelevant for a new target population.
\end{alertblock}
\end{frame}

\section{Populations and Estimands}

\begin{frame}{Study and Target Populations}
Let \(S=1\) indicate membership in the study population and \(S=0\) the target population.

The study estimand might be
\[
\E[Y(1)-Y(0)\mid S=1].
\]

The target estimand is
\[
\E[Y(1)-Y(0)\mid S=0].
\]
\end{frame}

\begin{frame}{The Estimand Must Be Named}
Transportability depends on what we want to transport:
\begin{itemize}
\item ATE in a new population;
\item ATT for a new treated group;
\item policy value under a deployment rule;
\item site-specific effect;
\item effect at a future time.
\end{itemize}
\end{frame}

\section{Effect Modification}

\begin{frame}{Effect Modifiers}
A covariate \(X\) is an effect modifier if
\[
\E[Y(1)-Y(0)\mid X=x]
\]
varies with \(x\).

If the distribution of \(X\) differs between study and target populations, the average treatment effect may change.
\end{frame}

\begin{frame}{Transport Requires Measured Modifiers}
If all relevant effect modifiers are measured, we can attempt to adjust for population differences.

If important effect modifiers are unmeasured, transport rests on stronger, less testable assumptions.
\end{frame}

\section{Standardization and Weighting}

\begin{frame}{Standardization}
Estimate conditional effects in the study population:
\[
\tau(x)=\E[Y(1)-Y(0)\mid X=x,S=1].
\]

Average them over the target population:
\[
\E_{X\mid S=0}[\tau(X)].
\]
\end{frame}

\begin{frame}{Transport Weights}
A weighting approach reweights study units to resemble the target population:
\[
w(X)=\frac{P(S=0\mid X)}{P(S=1\mid X)}.
\]

Then study effects are averaged with target-like covariate weights.
\end{frame}

\begin{frame}{Support Condition}
Transport requires that target covariate profiles have support in the study population.

If
\[
P(S=1\mid X=x)=0
\]
for some target-relevant \(x\), the study contains no empirical basis for that subgroup.
\end{frame}

\section{Covariate Shift and Site Heterogeneity}

\begin{frame}{Covariate Shift}
Covariate shift means
\[
P(X\mid S=1) \neq P(X\mid S=0).
\]

This matters for transport when treatment effects vary with \(X\).
\end{frame}

\begin{frame}{Site Heterogeneity}
Effects may differ across sites because of:
\begin{itemize}
\item population composition;
\item treatment implementation;
\item baseline risk;
\item institutions;
\item measurement systems;
\item competing interventions.
\end{itemize}
\end{frame}

\begin{frame}{Composition vs Mechanism}
A study and target may differ because:
\begin{itemize}
\item the mix of units differs;
\item the causal mechanism itself differs.
\end{itemize}

Standardization handles composition differences better than mechanism changes.
\end{frame}

\section{Trial-to-Target Generalization}

\begin{frame}{Trial Evidence for a Target Population}
Randomized trials can identify effects internally, but trial participants may differ from the deployment population.

Common differences:
\begin{itemize}
\item eligibility criteria;
\item adherence;
\item baseline risk;
\item access to care;
\item time period.
\end{itemize}
\end{frame}

\begin{frame}{Generalizability vs Transportability}
Generalizability often means the target population contains the study population.

Transportability often means evidence must be moved to a related but distinct population.

The mathematical issues are similar: measured effect modifiers and support dominate.
\end{frame}

\section{Temporal Transport}

\begin{frame}{Effects Can Drift Over Time}
A causal effect estimated in 2024 may not equal the effect in 2027.

Reasons include:
\begin{itemize}
\item technology changes;
\item competitor response;
\item policy adaptation;
\item population learning;
\item saturation effects.
\end{itemize}
\end{frame}

\begin{frame}{Stable Mechanisms}
Temporal transport is more plausible when the mechanism linking treatment to outcome remains stable.

It is less plausible when context changes the treatment itself or the response to it.
\end{frame}

\section{Implementation and Versions of Treatment}

\begin{frame}{Different Versions of Treatment}
The same treatment label may hide different implementations.

Examples:
\begin{itemize}
\item different marketing copy;
\item different clinicians;
\item different software interfaces;
\item different compliance regimes;
\item different pricing rules.
\end{itemize}
\end{frame}

\begin{frame}{Implementation Is Part of the Estimand}
If the treatment version changes, the transported effect may refer to a different intervention.

\begin{alertblock}{Practical warning}
Do not transport an effect across settings where the treatment label stayed fixed but the treatment mechanism changed.
\end{alertblock}
\end{frame}

\section{Sensitivity and Diagnostics}

\begin{frame}{Transport Diagnostics}
Useful checks include:
\begin{itemize}
\item compare study and target covariate distributions;
\item inspect overlap and transport weights;
\item identify target subgroups poorly represented in the study;
\item estimate site-specific effects when possible;
\item compare implementation conditions.
\end{itemize}
\end{frame}

\begin{frame}{Weight Instability}
Large transport weights indicate that a small number of study units represent large parts of the target population.

Consequences:
\begin{itemize}
\item high variance;
\item fragile conclusions;
\item extrapolation risk;
\item sensitivity to model specification.
\end{itemize}
\end{frame}

\begin{frame}{Unmeasured Effect Modifiers}
Sensitivity analysis can ask how strongly an unmeasured modifier would need to differ between study and target populations to change the transported conclusion.
\end{frame}

\section{Applied Reasoning}

\begin{frame}{Example: Regional Pricing Experiment}
Suppose a pricing experiment was run in urban regions.

Can we transport it to rural regions?

Ask:
\begin{itemize}
\item Are baseline demand and income different?
\item Are competitors different?
\item Is customer substitution different?
\item Is the price implementation identical?
\item Are rural profiles represented in the urban experiment?
\end{itemize}
\end{frame}

\begin{frame}{A Credible Transport Workflow}
\begin{enumerate}
\item Define the source population and target population.
\item Define the target estimand.
\item List plausible effect modifiers.
\item Compare modifier distributions.
\item Check support and transport weights.
\item Standardize or weight the study effect to the target.
\item Assess implementation differences.
\item Run sensitivity analysis for unmeasured modifiers.
\item Report where the evidence does not transport.
\end{enumerate}
\end{frame}

\begin{frame}{Common Failure Modes}
\begin{itemize}
\item assuming internal validity implies external validity;
\item transporting an ATE without specifying the target population;
\item ignoring effect modifiers;
\item extrapolating into unsupported target regions;
\item ignoring treatment-version differences;
\item treating old estimates as stable despite temporal drift;
\item hiding unstable transport weights.
\end{itemize}
\end{frame}

\begin{frame}{Synthesis: Effects Travel by Assumption}
\begin{itemize}
\item Internal validity and external validity are different problems.
\item Transport requires a named target estimand.
\item Effect modification is the central reason average effects differ across settings.
\item Standardization and weighting adjust for measured modifier distributions.
\item Support failures create extrapolation, not evidence.
\item Treatment implementation and time can change the causal mechanism.
\item External validity should be argued, diagnosed, and bounded, not assumed.
\end{itemize}
\end{frame}

\begin{frame}{Selected References}
\small
\begin{itemize}
\item Pearl, J., and Bareinboim, E. (2011). Transportability of Causal and Statistical Relations. \textit{AAAI}.
\item Bareinboim, E., and Pearl, J. (2016). Causal Inference and the Data-Fusion Problem. \textit{PNAS}, 113(27), 7345--7352.
\item Stuart, E. A. et al. (2011). Use of Propensity Scores to Assess the Generalizability of Results from Randomized Trials. \textit{JRSS A}, 174(2), 369--386.
\item Dahabreh, I. J. et al. (2020). Generalizing Causal Inferences from Individuals in Randomized Trials to All Trial-Eligible Individuals. \textit{Biometrics}, 76(2), 685--694.
\item Tipton, E. (2013). Improving Generalizations from Experiments Using Propensity Score Subclassification. \textit{JRSS A}, 176(2), 495--519.
\end{itemize}
\end{frame}

\contactslide

\end{document}
